\section{Affine test lattices and almost simple intervals}\label{sec:affine}

The negative examples will come from two opposite orientations of a finite
tree of truncated affine geometries. The common lattice is simple and
self-dual. Its reverse links have several horizontal branches, and their
bottom diamonds exclude particular normal-subgroup configurations. We first
construct these lattices and reduce a hypothetical representation to an
ordinary almost simple subgroup interval. Subsequent sections prune paths
through alternating and classical groups; bounded Lie rank gives a separate
contradiction, and additional path protections bound the number of remaining
stages.

\begin{theorem}\label{f:negative-family}
There are choices of the finite parameters in \Cref{f:def-test-lattice} for
which the resulting finite lattice is not the congruence lattice of any
finite algebra.
\end{theorem}

The construction and reductions in this section are uniform under the
inequalities stated below. The successive choices of larger parameters and
the completion of the proof of \Cref{f:negative-family} are given in
\Cref{sec:patterns,sec:classical,sec:alternating,sec:bounded,sec:weights}.
This theorem gives negative examples and, after bounding the parameters, a
cardinality ceiling. It does not assert that the chosen family member has
minimum cardinality; the minimum problem is treated separately.

\subsection{Structural inputs and conventions}

We use the classification of finite simple groups and its standard
structural consequences, including Schreier's Theorem, the automorphisms
and central covers of finite groups of Lie type, and the classification of
finite two-transitive groups. The algebraic-group inputs include
Lang--Steinberg, highest-weight theory, Steinberg's theorems, and the
algebraic Peter--Weyl filtration. The bounded-rank argument will use the
sufficiently-general subgroup Theorem of Larsen--Pink
\cite{LarsenPink2011}. To specify the scope of these inputs, we fix the
following conventions.

\begin{convention}[Lie-theoretic realizations]\label{f:conv-lie-inputs}
We use split root-system realizations over algebraic closures, with both
simply connected and adjoint root data, the Chevalley root and parabolic
structure, and the isogeny and Steinberg-map classification. Field maps,
pinned diagram maps, and, in the relevant special characteristics, the
special isogenies of bounded types are combined with inner automorphisms.
A Steinberg map of positive height supplies the finite fixed group.
For a simple adjoint group at sufficiently large parameter, or in the
high-rank cases used here, the associated simple finite group is the derived
fixed group. Outside the finitely many exceptional ranks and parameters for
simplicity and Schur multipliers, the simply connected fixed group is its
perfect universal central cover. Inner-diagonal, field, graph, and applicable
special actions induce the usual automorphisms of the simple projective
group. No uniform covering assertion is made at the excluded small cases.
For high-rank natural signed types, type $B$ is used only in odd
characteristic; in characteristic two we use the symplectic realization of
type $C$. These are the usual structure and automorphism conventions
\cite{GLS1994,Steinberg1968}.
\end{convention}

\begin{convention}[Representation-theoretic inputs]\label{f:conv-representation-inputs}
Over the algebraic closure in defining characteristic, we use rational
highest-weight theory: Weyl groups, Weyl and induced modules, weight convex
hulls, and the triangular and distribution-algebra properties of highest
weights. We also use the restricted-simple theorem for simply connected
groups and their Frobenius kernels, Steinberg's Tensor Product Theorem and
finite-group restriction and parametrization theorems on the simply
connected finite cover, and the two-sided algebraic Peter--Weyl good
filtration of the regular function algebra
\cite{Steinberg1963,Jantzen2003,Humphreys2006}.
The precise finite $P$-restricted convention at high rank is given in
\Cref{sec:weights}. These inputs concern rational modules or the specified
full finite covers; they do not assert that arbitrary almost simple subgroup
inclusions arise from algebraic-group inclusions.
\end{convention}

On the permutation side we use Jordan's prime-cycle Theorem: a primitive
group of degree $d$ containing a cycle of prime length at most $d-3$
contains $A_d$. We use Bochert's primitive-index bound, in particular an
exponential lower bound on the index in $S_d$ of a primitive group not
containing $A_d$ in sufficiently large degree, and Burnside's prime-degree
Theorem. The finite two- and multiply-transitive classification gives the
following consequences: a sufficiently large alternating simple group has
only its natural degree as a nontrivial two-transitive action; in high
degree a six-transitive group contains the alternating group; and a
non-affine two-transitive group of degree $37$ contains $A_{37}$
\cite{DixonMortimer1996}.
For the last assertion, the projective degrees $1+s+\cdots+s^j$, with $s$ a
prime power, never equal $37$: for $j=1$, one would need $s=36$; for
$j=2,3,4$ consecutive possible integer values bracket $37$ without equality;
and for $j\ge5$ even $s=2$ gives at least $63$. The unitary, Suzuki, Ree, and
symplectic-over-two degrees $s^3+1$, $s^2+1$, and
$2^{e-1}(2^e\pm1)$ also give no permitted parameter of degree $37$, and the
additional small and sporadic actions on the classification list do not
have that degree. The simple-group identifications in the classification
exclude spurious high-degree alternating overlaps.
We use ordinary finite subset Ramsey theory, finite multidimensional
Hales--Jewett for a fixed alphabet, number of colors, and number of free
slots, and later the finite vector-space Ramsey Theorem for linear
subspaces \cite{Ramsey1930,HalesJewett1963,GrahamLeebRothschild1972,GrahamLeebRothschild1973}.
More specialized uses of these inputs are stated where they occur.

\subsection{Charts, branches, and prepared links}

\begin{definition}[Affine test lattice]\label{f:def-test-lattice}
Let $q\ge5$ be a prime power such that $q-1$ is not a prime power, and let
$r,D,g,m,n,N$ be positive integers satisfying
\begin{equation}\label{f:eq-parameters}
 D\ge2,\qquad g\ge2,\qquad 3\le m<n-2,\qquad Dr<m,\qquad N\ge n+2.
\end{equation}
There are arbitrarily large allowable $q$, for example powers of $7$.
A \emph{chart} is a copy of $\operatorname{AG}(N,q)$. The rank of a nonempty
affine flat is its affine dimension plus one; the empty flat has rank zero.
To attach a fresh chart at a specified flat $W$, identify a flat of the new
chart affinely with $W$, and take all its other points fresh. Its intersection
with the union of the old point sets is exactly $W$. Attachment at the empty
flat makes no identifications.

Build the trunk charts recursively, beginning with the path ending at
$W=\varnothing$.
\begin{enumerate}[label=\textup{(\roman*)}]
\item At a path ending at $W$, attach $g$ fresh \emph{private charts} at $W$.
\item In each such chart, reserve a child path extending this path by $U$
for every flat $U\supset W$ with rank increment between $1$ and $r$.
Perform stages for $D$ generations: the empty path has a stage, as does
every child path reached in fewer than $D$ steps.
\item Finally, at every nonempty flat of rank at most $m$ in a trunk chart,
attach at least one additional fresh \emph{leaf chart}, with no stages of
its own.
\end{enumerate}
Repeated reservations may use distinct fresh charts. There are only finitely
many charts. The primal interior labels are the point sets that are
nonempty affine flats of dimension at most $n$ in some chart, ordered by
inclusion. Adjoin the empty bottom and one formal top above all these
labels, obtaining a bounded poset $L^+$. Let $L^-$ be an order-reversed copy
of $L^+$. The \emph{test lattice} $L$ is their horizontal sum: identify the
two bottoms and the two tops, with no other cross-comparabilities.
The interior labels on $L^-$ are called \emph{reverse flats}, named by their
affine point sets; their lattice order reverses inclusion. On this side the
top may accordingly be called the empty flat.
\end{definition}

A comparability component below means a connected component of the
undirected comparability graph on the strict interior of a bounded lattice.
Distinct such components are horizontal: two elements from different
components have meet zero and join one.

\begin{proposition}[Geometry of the test lattice]\label{f:prop-chart-properties}
The construction in \Cref{f:def-test-lattice} gives a finite self-dual
lattice $L$. It is graded of length $n+2$, atomistic, coatomistic, and simple.
Any two charts intersect in a flat of each, of rank at most $m$, with the
same affine structure on the intersection. Within one chart the primal
join of two flats is their affine span if that span has dimension at most
$n$. Attachment preserves all old meets and joins. In particular, flats in
two distinct private charts attached at the same $W$, neither contained in
$W$, have no common proper primal upper bound. There are at least three
horizontal comparability components in the interior of $L$, each containing
a chain of distinct comparable interior elements.
\end{proposition}
\begin{proof}
Induct on the attachments. When a chart is attached at $W$ in an old chart,
its intersection with any other old chart is exactly that chart's
intersection with $W$. By the inductive hypothesis this is a flat in the
old chart containing $W$ and in the other old chart, with agreeing affine
structures. It is therefore a flat of the new chart as well. This proves
the pairwise-intersection assertion. Trunk attachment flats have rank less
than $Dr<m$, and leaf attachment flats have rank at most $m$, so every
intersection of distinct charts has rank at most $m$. Affine subflats of a
shared flat are consequently identified consistently.

It follows that intersections of labels are again flats, allowing the
empty set. Moreover, whenever one nonempty label is contained in another,
it is an affine subflat in every chart containing the larger label. Thus
$L^+$ has finite meets and a top, which suffice to give all joins. A new
proper upper bound of two old labels, not itself already an old flat, can
contain those old labels only if both lie in the attachment flat $W$; it
then contains their old affine span in $W$. Hence no old join changes.
Meets are set intersections and do not change either. This also proves
the claim about private charts, first at their attachment and then after
all subsequent attachments. Within a chart the stated affine span is the
least proper upper bound. The horizontal sum with the opposite lattice
is a self-dual lattice.

Every proper chain in one orientation can be refined one affine dimension
at a time, up to dimension $n$, before reaching the formal top. Thus $L$
is graded of length $n+2$. A flat is generated by its points, so $L^+$ is
atomistic; intersecting its $n$-dimensional extensions inside one chart
recovers the flat, so it is coatomistic. The same properties hold in the
opposite orientation and hence in the horizontal sum.

The $g\ge2$ initial charts are disjoint. An attachment meets all old points
in a flat contained in one old chart; it cannot connect two old initial
branches by comparabilities. Each orientation therefore has at least two
horizontal interior components, and each component contains long affine
chains. In particular $L$ has at least three such components.

Finally, let a nonidentity lattice congruence identify two different
comparable elements $u<v$; such a pair exists whenever any two distinct
elements are identified, by meeting and joining them. By atomisticity there
is an atom $a\le v$ with $a\not\le u$. Meeting the identification with $a$
gives $a\equiv0$. Choose an interior element $b$ in another component.
Joining with $b$ gives $1\equiv b$. Meeting this with an interior element $c$
in a third component gives $c\equiv0$. Since $a\vee c=1$, we conclude
$1\equiv0$. The congruence is total, proving simplicity.
\end{proof}

\begin{proposition}[Prepared reverse links]\label{f:prop-prepared-links}
If $X$ is a nonempty flat of rank at most $m$ in a trunk chart, the reverse
ideal with top $X$, denoted $[0,X]$, is graded, atomistic, and coatomistic.
Its strict interior has at least two horizontal comparability components,
each containing comparable distinct elements. The following additional
properties hold for $L$ and for every such prepared link.
\begin{enumerate}[label=\textup{(\roman*)}]
\item\label{f:link-diamonds} Every bottom interval of length two is a diamond
$M_b$ with $b\ge3$ atoms, where $b-1\ge2$ is not a prime power.
\item\label{f:link-upper-diamonds} In every interior component there is an
element $y$ two ranks below the interval's top whose upper interval is a
diamond with at least two middle elements.
\end{enumerate}
In addition, $L$ has the following separation property:
\begin{enumerate}[label=\textup{(\roman*)},resume]
\item\label{f:link-separation} For any atom $u$ and any proper upper bound
$\omega\ge u$, there is a rank-two element $z\ge u$ with
$z\not\le\omega$.
\end{enumerate}
\end{proposition}
\begin{proof}
A leaf attached at $X$ supplies a branch of strict superflats of $X$
incomparable with the previous strict superflats. Later attachments cannot
connect two old branches in this reverse link: any connection through a
new chart passes through superflats of $X$ lying in its attachment flat;
if any exist, they are already connected within that flat. Every
superflat of $X$ lies in a chart containing $X$, and in that chart it has
extensions and intermediate flats. Each branch therefore contains
comparable distinct elements and chains out to $n$-flats. Gradedness follows
from these dimension chains. The flats extending $X$ by one dimension, and
the intersections of $n$-extensions, establish atomisticity and
coatomisticity in the reversed order.

On the primal side of $L$, a bottom interval of length two is a line, with
$b=q$ points; the choice of $q$ gives the claimed condition on $b-1$.
On the reverse side, including a prepared link, its top flat has dimension
$n-1$. Such a flat belongs to exactly one chart, since distinct charts meet
in rank at most $m<n-2$. Its $n$-dimensional extensions are counted by
one-dimensional subspaces of a vector space of dimension $N-n+1$, so
\[
 b=\frac{q^{N-n+1}-1}{q-1},\qquad
 b-1=q(1+q+\cdots+q^{N-n-1}).
\]
The second factor is greater than one because $N\ge n+2$, and is congruent
to $1$ modulo the characteristic prime. Thus $b-1$ has at least two prime
divisors. This proves \ref{f:link-diamonds}.

For \ref{f:link-upper-diamonds}, choose a primal $(n-1)$-flat, or, on the
reverse side, a flat whose rank exceeds that of the empty flat or the
prepared $X$ by two. The appropriate upper interval is a diamond, with at
least two middle flats. Such choices can be made in every branch.

For \ref{f:link-separation}, a primal atom is a point. In one chart through
it there is a line escaping the given proper flat $\omega$. A reverse atom
is an $n$-flat, belonging to a unique chart; its proper upper bound
corresponds to a nonempty subflat. Choose a hyperplane of the $n$-flat that
does not contain that subflat. It is the required rank-two reverse label.
\end{proof}

\subsection{From algebras to ordinary group intervals}

We write $[H,G]$ for the ordinary interval of all subgroups between finite
groups $H\le G$. The adjective \emph{ordinary} distinguishes such intervals
from invariant-subgroup intervals and from the extension posets introduced
below.

\begin{lemma}[Subdirect products of simple groups]\label{f:lem-subdirect}
A subdirect subgroup of a finite power of a nonabelian simple group is a
product of full, possibly automorphism-twisted, diagonal subgroups on
disjoint parts partitioning the coordinates. Containment between two such
subgroups means refinement of the tied parts, with agreement of the
diagonal identifications on each refined part. If a subgroup contains the
untwisted full diagonal, all its diagonal identifications are untwisted.
\end{lemma}
\begin{proof}
Induct on the number of coordinates. Project away the last coordinate.
The kernel in the last simple factor is normal in that factor, by
surjectivity, so is either the whole factor or trivial. In the first case
the last factor splits off. In the second, the subgroup is the graph of a
surjective homomorphism from the shorter subdirect product to the last
simple group. By induction the shorter group is a product of nonabelian
simple diagonal factors. The images of those factors are commuting normal
subgroups of a simple group; exactly one has nontrivial image, and on that
factor the map is an isomorphism. This ties the last coordinate to just
that part. The description of containment follows by projecting to pairs
of coordinates: a required tie must hold in the smaller subgroup with the
same identifying automorphism. Containing the untwisted full diagonal
forces each such automorphism to be the identity.
\end{proof}

\begin{proposition}\label{f:prop-ordinary-interval}
If the test lattice $L$ is representable by a finite algebra, it is
isomorphic to an ordinary subgroup interval of a finite group.
\end{proposition}
\begin{proof}
We begin with the least-carrier restriction argument of
\citet[Lemma~1]{PalfyPudlak1980}. The final diagonal construction below
will turn the resulting permutation action into an ordinary interval.
Choose a representing algebra with smallest possible carrier $A$. Let $F$
be the monoid of all maps $A\to A$ preserving its represented congruences.
The unary algebra $(A;F)$ has exactly those congruences: each represented
congruence is preserved by definition, while $F$ contains all one-variable
translations of the original operations, whose preservation forces an
equivalence to be an original congruence.

Let $e\in F$ be idempotent. On $eA$ use the unary operations
$ef|_{eA}$, $f\in F$. Restriction induces a surjective lattice homomorphism
\begin{equation}\label{f:eq-retraction}
 \operatorname{Con}(A;F)\longrightarrow
 \operatorname{Con}(eA;\{ef|_{eA}:f\in F\}).
\end{equation}
Meets are preserved by restriction. To check joins, retract an equivalence
path by $e$: its vertices then lie in $eA$ and each step remains in the
same congruence. For surjectivity, if $\delta$ is a congruence on the displayed
algebra on $eA$, define
\[
 x\mathrel{\widehat\delta}x'
 \quad\Longleftrightarrow\quad
 (efx,efx')\in\delta\ \text{for every }f\in F.
\]
This is an $F$-invariant equivalence. Its restriction to $eA$ is exactly
$\delta$: one direction uses its compatibility with all $ef|_{eA}$, and
the other uses $f=1$ and $e|_{eA}=1$.

If $|eA|>1$, the target of \eqref{f:eq-retraction} is nontrivial. Simplicity
of $L$ makes the homomorphism injective, and minimality of $|A|$ forces
equality of carrier sizes. A surjective idempotent is the identity.
Thus every nonunit idempotent of $F$ is constant.

Let $I$ be the ideal of nonunits of $F$. Every sufficiently long product
of elements of $I$ is constant. Indeed, otherwise take more than $|F|$
successive prefixes in order of application of a product whose final
image has size greater than one. Every prefix image then has size greater
than one, and two prefix maps coincide. The intervening nonempty product
fixes their common image pointwise. An idempotent power of that product
is a nonunit idempotent fixing at least two points, a contradiction.

For $\theta\in\operatorname{Con}(A;F)$ let $\Phi(\theta)$ be the equivalence
generated by
\[
 \{(f(x),f(y)):f\in I,\ (x,y)\in\theta\}.
\]
This is a congruence, by the ideal property of $I$, and
$\Phi(\theta)\le\theta$. It preserves joins: both equivalence generation and
images of paths show that applying it to a join gives the join of the
images. Iterating $\Phi$ generates precisely the pairs obtained from
products of the corresponding length in $I$. The inserted equivalence
closures contribute nothing extra, since an image of a path is a path;
diagonal pairs accommodate collapsed steps. The product observation makes
some iterate of $\Phi$ identically zero. For an atom $\alpha$, the deflation
$\Phi(\alpha)\le\alpha$ must therefore be zero, since the only other
possibility would fix $\alpha$ forever. Atomisticity now gives
$\Phi(1)=0$. Hence every map in $I$ sends every pair to a diagonal pair:
all nonunits are constant.

The units consequently form a permutation group $J$ with exactly $L$ as
its lattice of invariant equivalences; constants impose no conditions.
Choose any finite nonabelian simple group $C$, and form the coordinate
permutation semidirect product $C^A\rtimes J$. Then
\begin{equation}\label{f:eq-diagonal-interval}
 [\Delta(C)\rtimes J,\ C^A\rtimes J]
\end{equation}
is dual to that invariant-equivalence lattice. Indeed, intersection with
$C^A$ identifies this interval with the $J$-invariant subgroups above the
full diagonal. By \Cref{f:lem-subdirect} these are the diagonal products
determined by partitions of $A$, and subgroup containment reverses
refinement of the corresponding equivalences. Since $L$ is self-dual,
\eqref{f:eq-diagonal-interval} represents $L$ itself.
\end{proof}

\subsection{The monolith and extension description}

For a subgroup $H\le G$, its core is
$\operatorname{core}_G(H)=\bigcap_{g\in G}gHg^{-1}$. A group with a unique
minimal nontrivial normal subgroup is called \emph{monolithic}; that
minimal normal subgroup is its \emph{monolith}.
The normal-supplement and unique-monolith argument follows
\citet[Lemma~3.6]{PalfyTwisted}. The horizontal components of our prepared
intervals give the required lattice obstruction directly.

\begin{proposition}[Normal supplements]\label{f:prop-monolith}
Suppose $[H,G]$ is isomorphic either to $L$ or to a prepared reverse link.
After quotienting by $\operatorname{core}_G(H)$, every nontrivial normal
subgroup is supplemented by $H$. The quotient has a unique minimal normal
subgroup $M=S^s$, with $S$ nonabelian simple, and $H$ acts transitively on
its simple factors. Moreover $C_G(M)=1$ in the core quotient. In that
quotient the subgroup interval identifies with
\begin{equation}\label{f:eq-invariant-interval}
 [P,S^s]^H,\qquad P=H\cap S^s,
\end{equation}
the $H$-invariant subgroups in this range, by intersection with $S^s$;
the inverse sends $Y$ to $HY$.
\end{proposition}
\begin{proof}
If $R\trianglelefteq G$, then $HR$ permutes as a set with every overgroup of
$H$. Dedekind's law consequently makes $HR$ left-modular in $[H,G]$.
There is no proper interior left-modular label: for a putative such label
$u$, choose comparable distinct interior labels $v<w$ in another horizontal
component. Then
\[
 (v\vee u)\wedge w=w\ne v=v\vee(u\wedge w).
\]
Thus either $R\le H$ or $HR=G$. In the core quotient the first alternative
forces $R=1$.

An abelian minimal normal supplement would identify the interval with
an invariant-subgroup interval in that abelian group. Such an interval is
modular, contradicting the displayed failure of the modular identity.
If two distinct minimal normal subgroups existed, they would commute.
Using one to supplement $H$, the conjugation action of $H$ on the other
would induce all its inner actions. The interval would again identify
with a modular interval, now of invariant normal subgroups of that other
minimal normal subgroup. Both possibilities are impossible. The unique
minimal normal subgroup is therefore a power $S^s$ of a nonabelian simple
group. Its factor orbits under $H$ are normal in $HM=G$, so minimality makes
the action transitive. Its centralizer is normal; if nontrivial it would
contain $M$, which is impossible because $M$ is centerless. Finally $HM=G$,
and the elementary intersection--product correspondence gives
\eqref{f:eq-invariant-interval}.
\end{proof}

For a subgroup $Y\le S^s$, write $c(Y)$ for the product of its coordinate
projections. It is $H$-invariant whenever $Y$ is. We will need a precise
parametrization of the invariant subdirect groups. We use the
invariant-subdirect/extension correspondence in
\citet[Proposition~2.2 and Lemma~2.3]{PalfyTwisted}, expressed here for
the transitive action on the actual simple coordinates.

\begin{lemma}[Extension description]\label{f:lem-extension-description}
Suppose $H$ acts by automorphisms on $S^s$, transitively on the simple
coordinates. Fix a coordinate $o$, identify its group with $S$, and let
$A_0\le H$ be its stabilizer, with induced action
$\alpha:A_0\to\operatorname{Aut}(S)$. The $H$-invariant subdirect subgroups
of $S^s$ are in bijection with extensions
\begin{equation}\label{f:eq-extension-data}
 A_0\le T\le H,\qquad
 \beta:T\longrightarrow\operatorname{Aut}(S),\qquad
 \beta|_{A_0}=\alpha.
\end{equation}
This bijection reverses extension order: increasing the extension means
increasing its domain and agreeing on the smaller domain. If one restricts
to subdirect groups containing a specified $P$, then above any retained
group in actual subgroup order all restrictions of its extension to
intermediate domains still occur.
\end{lemma}
\begin{proof}
By \Cref{f:lem-subdirect}, a subdirect group has tied diagonal parts that
form an $H$-invariant partition of the coordinates. The part containing $o$
is a block for the transitive action; it equals $To$, where $T$ is its
setwise stabilizer and $A_0\le T$. The action of $T$ on that diagonal,
identified with $S$ by projection to $o$, gives $\beta$ extending $\alpha$.

Conversely, given \eqref{f:eq-extension-data}, define the diagonal on $To$
using left-action notation. At coordinate $ho$, for $h\in T$, send $u\in S$
to the transport by $h$ of $\beta(h)^{-1}(u)$ from coordinate $o$. Replacing
$h$ by $ha$, with $a\in A_0$, gives the same element because
$\beta(a)=\alpha(a)$. This defines a $T$-invariant full diagonal. Transport
it to the other blocks and multiply the resulting disjoint diagonal
factors. The result is $H$-invariant and subdirect, and recovers the original
data. Refinement of tied parts, with agreement of identifications, is
exactly restriction of the extension. The last assertion follows because
such restrictions give larger actual subgroups, still containing $P$.
\end{proof}

\begin{proposition}[Cases for a prepared interval]\label{f:prop-interval-cases}
Use the notation of \Cref{f:prop-monolith}, in the core quotient. Let
$\Lambda\le\operatorname{Aut}(S)$ be the image of a coordinate stabilizer
in $H$, and put $I=\Lambda\cap S$, identifying $S$ with its inner
automorphisms. The following are the only possibilities.
\begin{enumerate}[label=\textup{(\roman*)}]
\item\label{f:case-subdirect} $c(P)=S^s$. All labels are subdirect. If the
interval represents the whole $L$, its opposite has an ordinary subgroup
representation with top-group order at most $|H|<|G|$.
\item\label{f:case-product} $P=c(P)\ne1$. The interval is isomorphic to the
ordinary almost simple interval $[\Lambda,S\Lambda]$.
\item\label{f:case-extension} $P=1$ and $I=S$. Every nonbottom label is
subdirect, and the dual interval is the extension poset of
\Cref{f:lem-extension-description} with one artificial top adjoined.
\end{enumerate}
In particular the cases $P<c(P)<S^s$, $P=1$ with $I=1$, and $P=1$ with
$1<I<S$ are impossible.
\end{proposition}
\begin{proof}
We check all the cases, retaining the prepared-interval hypotheses.

\emph{Subdirect bottom.}
If $c(P)=S^s$, every label is subdirect and refines the fixed diagonal ties
of $P$. Restrict to one tied part of $P$. Its setwise stabilizer in $H$
acts transitively on its coordinates. A refinement of its coordinates
that is invariant under this stabilizer extends uniquely to the other
parts by $H$, keeping the twists fixed by $P$. Thus the interval is dual
to the invariant-equivalence lattice of that transitive action. By
\Cref{f:lem-coset}, the latter is an ordinary interval above a point
stabilizer in the part stabilizer. Its top has order at most $|H|<|G|$,
proving \ref{f:case-subdirect}.

\emph{A proper nonproduct bottom.}
Suppose $P<c(P)<S^s$. Choose a horizontal interior component different from
that of $c(P)$. Every label $Y$ in it has
$Y\vee c(P)=S^s$ and $Y\cap c(P)=P$. Its coordinate projections already
contain those of $P$, so adjoining $c(P)$ does not enlarge them. The join
condition therefore makes $Y$ subdirect. In its diagonal description,
intersection with $c(P)$ restricts each coordinate to the corresponding
nontrivial projection of $P$, with the same ties. These restrictions
project onto every such subgroup, because the intersection contains $P$.
There are comparable distinct labels in the chosen component. Their
strict containment changes at least one diagonal tie, and releasing a tie
strictly enlarges the intersection: its tied coordinate subgroup is
nontrivial and may now be chosen independently on the separated parts.
This contradicts the fact that both intersections equal $P$.

\emph{A nontrivial product bottom.}
Suppose $P=c(P)\ne1$, necessarily a proper product. A subdirect group
containing independent nontrivial subgroups in each coordinate cannot
have any nontrivial ties. Thus the only subdirect label above $P$ is
$S^s$. For a coatom $Y$, the product $c(Y)$ is a label containing $Y$;
maximality and the preceding observation imply $Y=c(Y)$. Since the
interval is coatomistic, every label is a product.
At one coordinate let $Q$ be the projection of $P$. Then
$1<Q\le I$: the coordinate elements of $P\le H$ induce their inner
automorphisms. Products identify the interval with $[Q,S]^\Lambda$.
For any coatom $T$ of this invariant interval, $I$ normalizes $T$.
If $TI=S$, then $T$ is normal in $TI$, hence in $S$, impossible because
$1<Q\le T<S$. Thus $TI<S$, and maximality gives $I\le T$.
Intersecting all coatoms yields $I\le Q$, so $I=Q$.
Intersection with $S$ and multiplication by $\Lambda$ now identify
$[Q,S]^\Lambda$ with $[\Lambda,S\Lambda]$ inside
$\operatorname{Aut}(S)$, proving \ref{f:case-product}.

\emph{Trivial bottom and no inner operators.}
Suppose $P=1$ and $I=1$. Schreier's Theorem makes $\Lambda$ soluble, since
it embeds in $\operatorname{Out}(S)$. We first show that a soluble operator
group on a finite nonabelian simple group has a nontrivial proper invariant
subgroup. We may replace it by its effective image. If that image is
trivial, choose any subgroup of prime order. Otherwise choose a minimal
normal elementary abelian $p$-subgroup $A$ of the operator group. Its fixed
group is invariant, and if nontrivial it is proper by effectiveness.
If its fixed group is trivial, orbit counting on the simple group gives
$p\nmid|S|$. For each prime divisor $\ell$ of $|S|$, the number of Sylow
$\ell$-subgroups divides $|S|$ and is prime to $p$, so the $p$-group $A$
fixes at least one such Sylow subgroup. There is at most one: the set of
conjugators between two $A$-invariant Sylow subgroups is $A$-invariant and
has size $|N_S(P_\ell)|$, prime to $p$, so has an $A$-fixed conjugator.
The trivial fixed group forces that conjugator to be $1$. Uniqueness
makes the Sylow subgroup invariant under the whole operator group.
This proves the assertion.

Choose a minimal nontrivial $\Lambda$-invariant subgroup $Q$ of $S$.
Its characteristic subgroups are invariant, so $Q$ is characteristically
simple. If it were a power of a nonabelian simple group, apply the
preceding assertion to the soluble operator image of a factor stabilizer,
and transport the resulting proper invariant subgroup through the
transitive factor action. This would give a nontrivial proper invariant
subgroup of $Q$, a contradiction. Therefore $Q$ is elementary abelian.
Its normalizer is proper because $S$ is simple. It strictly contains $Q$:
if $Q$ is not Sylow, use the normalizer property in a Sylow subgroup
containing it; if it is Sylow and were self-normalizing, Burnside transfer
would give a normal prime complement, again contradicting simplicity.

Transport $Q$ and $N_S(Q)$ to all coordinates and take their products.
They are labels, the latter strictly above the former. If the product of
the $Q$'s has lattice rank at least two, choose a rank-two label below it
and any atom in that interval. Otherwise the product itself is an atom;
choose a rank-two label above it and below the normalizer product. In both
cases the chosen atom $A$ is an elementary abelian normal subgroup of the
rank-two group $R$. Every other atom is an $H$-invariant complement to
$A$ in $R$, and conversely every such complement is an atom: intersection
with $A$ is trivial, and the join with $A$ is $R$ in the diamond.
Fix one invariant complement $C$. In the semidirect product $R=A\rtimes C$,
other complements are graphs of derivations $C\to A$. Since $H$ preserves
both factors, the invariant complements correspond precisely to the
$H$-equivariant derivations. These form a vector space over the prime field
of $A$. Their number is a prime power, allowing the zeroth power.
This is the operator-invariant version of the complement-counting
argument in \citet[Theorem~3]{PalfyPudlak1980}.
The number equals $b-1$ in the bottom diamond $M_b$, contradicting
\Cref{f:prop-prepared-links}\ref{f:link-diamonds}.

\emph{Trivial bottom and proper nontrivial inner operators.}
Suppose $P=1$ and $1<I<S$. Transporting $I$ gives a proper interior product
label. A label $Y$ in another horizontal component is subdirect: its
nontrivial coordinate projection $R$ satisfies $RI=S$, because its join
with the product of the $I$'s is full. The subgroup $I$ normalizes $R$,
so $R\trianglelefteq RI=S$ and hence $R=S$.
Write $(T,\beta)$ for the extension describing $Y$ in
\Cref{f:lem-extension-description}. Its image contains all inner
automorphisms. Otherwise $\beta(T)\cap S$ is a nontrivial proper invariant
subgroup of the diagonal, containing $I$. Transporting it over the tied
parts and multiplying gives a nonbottom label below $Y$ with proper
coordinate projections. It would lie in the same component, contradicting
the preceding conclusion about all its labels.

Choose $y$ in this component two below the top with a diamond above it.
The filter above $y$ identifies dually, by restriction of $\beta_y$, with
$[A_0,T_y]$. Every domain strictly larger than $A_0$ has image containing
$S$. Apply this observation to $A_0\ker\beta_y$: its image is still
$\Lambda$, which does not contain $S$. Hence $\ker\beta_y\le A_0$.
After quotienting by this kernel, every proper increase over $\Lambda$
contains $S\Lambda$, so there is a unique minimal proper increase. This
contradicts the diamond's having at least two middle elements.

\emph{Trivial bottom and all inner operators.}
Finally, if $P=1$ and $I=S$, every nonbottom label has nontrivial coordinate
projections invariant under all inners, so those projections equal $S$.
The nonbottom labels are exactly the invariant subdirect groups, and
\Cref{f:lem-extension-description} identifies their dual with the extension
poset. The bottom contributes the artificial top. This proves
\ref{f:case-extension} and completes the case analysis.
\end{proof}

Choose now an ordinary representation of the whole $L$ with smallest
possible top-group order, which exists under the hypothesis of
\Cref{f:prop-ordinary-interval}. Self-duality and
\Cref{f:prop-interval-cases}\ref{f:case-subdirect} exclude the subdirect-bottom
case. Consequently either $L$ already has an ordinary almost simple
representation, or it is isomorphic to an extension poset with an artificial
top. We eliminate the latter alternative unless a prepared reverse point
link has the desired almost simple representation.

\subsection{Pure extensions and their kernels}

The full-inner extension case and its kernel description appear in
\citet[Proposition~2.5 and Remark~2.6]{PalfyTwisted}; see also the
attribution to Aschbacher's signalizer formulation there in Section~4.
We now use the affine lattice to constrain these extensions further.

\begin{definition}\label{f:def-pure-extension}
An extension datum consists of a subgroup $H$ of a finite ambient group
$K$, a nonabelian simple group $B$, and a homomorphism
$\psi:H\to\operatorname{Aut}(B)$ whose image contains $B$, identified with
its inner automorphisms. Its extension poset consists of all homomorphisms
$\beta:T\to\operatorname{Aut}(B)$, $H\le T\le K$, restricting to $\psi$,
ordered by extension, together with an artificial top. Put
\[
 D_0=\ker\psi,\qquad P_0=\psi^{-1}(B).
\]
An extension is \emph{pure} if $\beta(T)=\psi(H)$, equivalently
$T=H\ker\beta$.
\end{definition}

\begin{lemma}[Purity and kernel joins]\label{f:lem-purity}
If the full extension poset with artificial top of
\Cref{f:def-pure-extension} is isomorphic to $L$, every proper label is pure.
These labels are then parametrized by all $H$-normalized subgroups $E\le K$
satisfying
\begin{equation}\label{f:eq-pure-kernel}
 D_0\le E,\qquad E\cap H=D_0.
\end{equation}
The domain is $T=HE$, and its entire lower ideal is the ordinary interval
$[H,T]$. Proper joins of kernel labels are generated subgroups. If two
kernel labels join to the artificial top, their generated subgroup
contains $P_0$.
\end{lemma}
\begin{proof}
Subextensions of pure extensions are pure. If finitely many pure labels
have a common extension, their join is its restriction to the subgroup
generated by their domains, and is pure as well: its image is generated
by copies of $\psi(H)$.

Suppose an atom $(T,\beta)$ is not pure. The intermediate restriction to
$H\ker\beta$ is pure, so atomhood forces $H\ker\beta=H$ and
$\ker\beta=D_0$. Since $H$ already maps onto the inner subgroup $B$, it
also follows that $\beta^{-1}(B)=P_0$. Thus $T$ normalizes both $P_0$ and
$D_0$. The simultaneous normalizer
\[
 N_K(P_0)\cap N_K(D_0)
\]
acts on $P_0/D_0\cong B$. Its conjugation homomorphism restricts to $\psi$
on $H$, since action on the identified inner automorphisms determines an
automorphism of the centerless group $B$. On each nonpure atom it similarly
agrees with $\beta$. This gives one extension label $\omega$, strictly below
the artificial top, bounding every nonpure atom.

Choose a nonpure atom $u$. By
\Cref{f:prop-prepared-links}\ref{f:link-separation}, there is a rank-two
label $z\ge u$ not below $\omega$. At most one atom of its diamond lies
below $\omega$, since two distinct atoms join to $z$. The diamond has at
least three atoms, so at least two are pure. Their proper join is $z$,
which is consequently pure; its subextension $u$ is pure, a contradiction.
All atoms, and hence all proper labels by atomisticity and purity of proper
joins, are pure.

A pure extension has a kernel satisfying \eqref{f:eq-pure-kernel} and
domain $HE$. Conversely, such a subgroup $E$ is normal in $HE$, and
$HE/E\cong H/D_0$ defines the unique extension killing $E$ and agreeing with
$\psi$. This proves the parametrization. If $H\le U\le HE$, then
$U=H(U\cap E)$, so restricting the extension identifies its lower ideal
with all of $[H,HE]$.
For a proper join, restriction of a common extension to the generated
domain shows that the kernel is the subgroup generated by the kernels:
their product with $H$ is the generated domain, and their intersection
with $H$ is $D_0$. If two labels instead join to the artificial top and
their generated kernel $E$ still met $H$ in $D_0$, the same construction
would give a common proper extension, impossible. Therefore $E\cap H$
properly contains $D_0$. It is normal in $H$, so its image is a nontrivial
normal subgroup of the almost simple group $\psi(H)$; it contains $B$.
As $D_0\le E$, this implies $P_0\le E$.
\end{proof}

\subsection{The lower affine-grid contradiction}

We compare pure extensions indexed by the points, lines, and planes of
an affine four-space. Four dimensions allow two planes to meet in exactly
one point; after using this to choose a suitable plane, the contradiction
is obtained inside that plane. The next lemma requires purity only for
the labels of dimension at most two.

\begin{lemma}[Lower-grid obstruction]\label{f:lem-lower-grid}
Fix an extension datum as in \Cref{f:def-pure-extension}, and an affine
four-space over a finite field with at least three elements. Suppose each nonempty
flat $X$ of dimension at most two supplies a pure extension with domain
$T_X>H$ and kernel $E_X$, ordered in reverse containment, with the following
properties.
\begin{enumerate}[label=\textup{(\roman*)}]
\item If $X\cap Y=\{i\}$, their labels have proper join the point label $i$;
thus $T_i=\langle T_X,T_Y\rangle$ and
$E_i=\langle E_X,E_Y\rangle$. Distinct point labels have incompatible
join, and their generated kernels contain $P_0$.
\item For each $X$, modulo $L_X=\operatorname{core}_{T_X}(H)$ the interval
has the normal-supplement and unique nonabelian monolith properties of
\Cref{f:prop-monolith}, including transitivity of the base on the simple
factors and the trivial monolith centralizer. Denote
its monolith by $W_X$, with preimage $V_X$ in $T_X$.
\item For every point $i\in X$, the subgroup
$(T_X\cap V_i)/L_i$ is subdirect in $W_i$ and is strictly above the
intersection of the base with $W_i$.
\end{enumerate}
These data cannot exist.
\end{lemma}
\begin{proof}
We organize the argument into the compatible monolith embeddings, the
selection of a plane with small cores, and the final factor action.

\emph{Monolith embeddings.}
For $i\in X$ put
\begin{equation}\label{f:eq-zxi}
 Z_{X,i}=(T_X\cap V_i)/L_i\le W_i.
\end{equation}
The base $H$ is transitive on the simple factors of $W_i$. It is therefore
transitive on the tied diagonal factors of its invariant subdirect group
$Z_{X,i}$. This makes $Z_{X,i}$ a nontrivial minimal normal subgroup of
\[
 T_X/L_i=(H/L_i)Z_{X,i}.
\]
It is not contained in $H/L_i$, by the strictness in the hypothesis.
Since $T_X\le T_i$, we have $L_i\le L_X$. Minimal normality and
$L_X\le H$ now give $L_X\cap V_i=L_i$. Thus the image of $Z_{X,i}$ in
$T_X/L_X$ is still injective and minimal normal, so equals $W_X$.
We obtain a canonical $H$-equivariant subdirect injection
$f_{X,i}:W_X\to W_i$, and
\begin{equation}\label{f:eq-monolith-preimages}
 V_X=(T_X\cap V_i)L_X.
\end{equation}

For $i\in X\subseteq Y$ we have
$H\le T_Y\le T_X\le T_i$ and $L_i\le L_X\le L_Y$.
Intersecting \eqref{f:eq-monolith-preimages} with $T_Y$ gives
\[
 V_X\cap T_Y=(V_i\cap T_Y)L_X,\qquad
 V_Y=(V_X\cap T_Y)L_Y,\qquad
 L_Y\cap V_X=L_X.
\]
The last equality follows also by using $L_Y\cap V_i=L_i$.
These formulas give injections $f_{Y,X}:W_Y\to W_X$ composing with the
$f_{X,i}$ as expected: lift representatives in $V_i\cap T_Y$.
They are subdirect. Indeed, every simple factor of $W_X$ is visible by
projection to some simple coordinate of $W_i$ through $f_{X,i}$, and the
composite image of $W_Y$ projects onto every such coordinate. The kernel
of the action of $H$ on $W_X$ is exactly $L_X$, by the trivial centralizer
of the monolith in $T_X/L_X$.

If $X\cap Y=\{i\}$, the proper-join condition gives
$T_i=\langle T_X,T_Y\rangle$. Writing each domain as $H$ times its
intersection with $V_i$ shows that
\begin{equation}\label{f:eq-core-intersection}
 W_i=\langle Z_{X,i},Z_{Y,i}\rangle,
 \qquad L_i=L_X\cap L_Y.
\end{equation}
For completeness, the subgroup generated by the two $Z$'s contains
$(H\cap V_i)/L_i$, is normalized by $H/L_i$, and multiplying it by $H/L_i$
gives $T_i/L_i$. Intersecting that product with $W_i$ therefore recovers
all of $W_i$. An element of $H$ acts trivially on this generated group
exactly when it acts trivially on both embedded monoliths; this proves
the asserted equality of cores.

\emph{A plane whose cores are killed by $\psi$.}
For any affine plane $U$, write
\[
 L_* =\bigcap_{i\in U}L_i.
\]
Every line domain $T_Y$, for a line $Y\subset U$, normalizes $L_*$. To see
this, fix a target point $j\notin Y$ and choose two distinct lines through
$j$ meeting $Y$, at points $k_1,k_2$. An element of $T_Y\le T_{k_t}$
normalizes $L_{k_t}$, so sends $L_*\le L_{k_t}$ into the core of the line
$jk_t$, because $L_{k_t}\le L_{jk_t}$. Equation
\eqref{f:eq-core-intersection} for those two lines puts the resulting
image inside $L_j$. For $j\in Y$, use $T_Y\le T_j$ directly. The same
argument for inverses proves normalization. Two distinct lines through a
point generate its domain, so $L_*$ is normal in the group generated by
all point domains of $U$.

Since $L_*\le H$, we have $L_*\cap E_i=L_*\cap D_0$ for every point $i$.
Both $E_i$ and $L_*$ are normal in $T_i$, so
$[E_i,L_*]\le L_*\cap D_0$. Thus every $E_i$ centralizes
$L_* /(L_*\cap D_0)$. The group generated by kernels of distinct points
contains $P_0$, by hypothesis. Therefore $\psi(L_*)$ centralizes $B$,
whose centralizer in $\operatorname{Aut}(B)$ is trivial. Hence
\begin{equation}\label{f:eq-plane-intersection-killed}
 L_*\le D_0.
\end{equation}

There is a point $i$ with $L_i\le D_0$. Otherwise every $\psi(L_i)$ is a
nontrivial normal subgroup of $\psi(H)$ and contains $B$. Take an iterated
commutator of the original normal subgroups $L_i$ inside $H$. It lies in
$\bigcap_{i\in U}L_i$ and its $\psi$-image still contains $B$, because
$B$ is perfect.
That contradicts \eqref{f:eq-plane-intersection-killed}.
Choose two affine planes in the four-space meeting exactly at this point
$i$. Their cores intersect in $L_i$ by
\eqref{f:eq-core-intersection}, so their two core images cannot both
contain $B$, by the same commutator argument. Hence one of these planes,
henceforth denoted $U$, has $L_U\le D_0$. All point cores of $U$ then lie
in $D_0$, since $L_i\le L_U$ for $i\in U$.

Quotient the group generated by the point domains of $U$ by
$\bigcap_{i\in U}L_i$. This subgroup is normal there as shown above,
lies in the base, and is contained in all the point cores and all extension
kernels in the plane. The core in each quotient is the quotient of the
old core. Thus the inclusions, join statements, and monolith identifications
survive. We retain the notation and may assume
\begin{equation}\label{f:eq-zero-common-core}
 \bigcap_{i\in U}L_i=1,\qquad L_U\le D_0.
\end{equation}

\emph{Compatible perfect subgroups.}
For a nonempty affine subflat $X\subseteq U$ set
\begin{equation}\label{f:eq-mx}
 M_X=E_X\cap\bigcap_{i\in X}V_i.
\end{equation}
Every factor $V_i\cap T_X$ is normal in $T_X$, as is $E_X$, so $M_X$ is
normal in $T_X$. It lies in $V_X$ by
\eqref{f:eq-monolith-preimages}. It maps onto $W_X$: each
$V_i\cap T_X$ maps onto $W_X$, while the normal image of $E_X$ modulo $L_X$
is nontrivial, since $HE_X=T_X>H$ and $L_X\le H$, and therefore contains
the monolith. Taking nested commutators of these finitely many normal
subgroups gives a subgroup in their intersection whose image still
contains $W_X$, by perfection of $W_X$. As $M_X\le V_X$, its image is exactly
$W_X$. The supplement condition then yields
\begin{equation}\label{f:eq-domain-mx}
 HM_X=T_X,\qquad D_0M_X=E_X.
\end{equation}
The second equality follows by intersecting the first with $E_X$.

For $X=U$, the kernel of $M_U\to W_U$ is contained in every $L_i$: it lies
in $L_U\cap V_i=L_i$. Equation \eqref{f:eq-zero-common-core} makes it trivial.
Consequently $M_U\cong W_U$ is an embedded product of nonabelian simple
groups; write its factors as $C_a$. The normal subgroups $L_U$ and $M_U$
of $T_U$ have trivial intersection, so $L_U$ centralizes every $C_a$.

The subdirect embedding $f_{U,X}$ assigns to each factor $C_a$ a disjoint
set of factors of $W_X$ on which it has its tied diagonal, and these sets
partition all factors of $W_X$. Let $M_{X,a}$ be the preimage in $M_X$ of
the entire product of $W_X$-factors assigned to $a$. These groups generate
$M_X$, are permuted by $H$, and each is normalized by $L_U$: the latter
centralizes $M_U$, so fixes its factors and their assigned sets.
For $i\in X$, the images of $M_{X,a}$ in $W_i$ have support only over the
same index $a$. For $a\ne b$ it follows that
\begin{equation}\label{f:eq-cross-commutators}
 [M_{X,a},C_b]\le L_i\le L_U.
\end{equation}
In fact $M_{X,a}$ centralizes $C_b$. Since $L_U$ centralizes $C_b$,
\eqref{f:eq-cross-commutators} and the three-subgroups lemma imply
$[M_{X,a},[C_b,C_b]]=1$; perfection gives the claim. Likewise, for meeting
lines $X,Y$ (allowing $X=Y$) and $a\ne b$, a common point $i$ gives
\begin{equation}\label{f:eq-line-commutators}
 [M_{X,a},M_{Y,b}]\le L_i\le L_U.
\end{equation}

\emph{The final action on the simple factors.}
Choose four pairwise nonparallel lines in $U$, two through a point $i$ and
two through a distinct point $j$. There are at least four directions,
since the field has at least three elements; the lines can be chosen
with the required distinct directions. Set
\[
 R'=\langle M_Y:Y\text{ is one of these four lines}\rangle.
\]
It is normalized by $H$. The two proper point joins and then the
incompatible join of $i,j$, together with \eqref{f:eq-domain-mx}, imply
$P_0\le D_0R'$. Hence
\begin{equation}\label{f:eq-surjects-b}
 \psi(P_0\cap R')=B.
\end{equation}

It remains to show that this image is soluble. We first show that
$P_0\cap R'$ fixes every factor $C_a$. Schreier's Theorem then controls
its outer action, and we show that $\psi$ kills the subgroup inducing
inner automorphisms.

Fix an index $b$.
A word in the generating groups $M_{Y,a}$ can be reordered as
\[
 \ell u v,
\]
where $\ell\in L_U$, all letters of $u$ have indices $a\ne b$, and all
letters of $v$ have index $b$. For adjacent letters in the wrong order,
put $e=xyx^{-1}y^{-1}\in L_U$, by
\eqref{f:eq-line-commutators}. If $p$ is the preceding word, then
\[
 pxy=peyx=e(e^{-1}pe)yx.
\]
Conjugating the preceding letters by $e$ preserves their indices because
$L_U$ normalizes every indexed generating subgroup. Thus the swap reduces
the inversion count. Induction gives the required reordering.

In right-conjugation notation, $C_b^{\ell u v}=C_b^v$, since $\ell$ and $u$
centralize $C_b$. The element $v$ centralizes every other $C_a$, and $C_b$
already commutes with them, so $C_b^v$ commutes with every $C_a$ for
$a\ne b$. If the word lies in $P_0\cap R'\le H$, it permutes the factors
of $M_U$; its image of $C_b$ cannot be another $C_a$, because that would
make a nonabelian simple group commute with itself. It must therefore be
$C_b$, proving the assertion.

The action of $P_0\cap R'$ on $W_U\cong M_U$ consequently has soluble
quotient modulo its inner action, by Schreier's Theorem on each fixed
simple factor. An element inducing an inner action on $W_U$ lies in
$V_U$: after multiplying by the corresponding monolith element its image
in $T_U/L_U$ centralizes $W_U$, and that centralizer is trivial. On the
other hand, purity and $L_U\le D_0$ give
\[
 H\cap V_U\le H\cap(E_UL_U)\le D_0L_U=D_0,
\]
because the nontrivial normal image of $E_U$ contains $W_U$ and hence
$V_U\le E_UL_U$. Thus the inner-action kernel within $P_0\cap R'$ is
killed by $\psi$. Consequently $\psi(P_0\cap R')$ is a quotient of a soluble group,
contradicting \eqref{f:eq-surjects-b} and the nonabelian simplicity of $B$.
\end{proof}

\begin{theorem}[Ordinary almost simple starting interval]\label{f:thm-almost-simple-reduction}
If the test lattice $L$ is representable, then either the whole $L$ or a
prepared reverse link with top an initial trunk point has an ordinary
almost simple subgroup-interval representation. On the reverse side,
there are at least $D-1$ further generations of prepared stages starting
at that top. This conclusion uses only the field condition and
\eqref{f:eq-parameters}, independently of later choices of larger
parameters.
\end{theorem}
\begin{proof}
Choose the least-top-order ordinary representation furnished by
\Cref{f:prop-ordinary-interval}. The reduction after
\Cref{f:prop-interval-cases} either already gives the whole $L$ as an
ordinary almost simple interval, or gives an extension poset with
artificial top isomorphic to $L$. In the latter case all proper labels
are pure by \Cref{f:lem-purity}. For each nonempty reverse flat $X$, write
$E_X,T_X$ for its kernel and domain. Its lower ideal is the ordinary
interval $[H,T_X]$.

Take an affine four-subspace of one initial trunk chart and consider all
its points, lines, and planes. Their ranks are at most three, hence at
most $m$, so all their reverse links are prepared. The proper and
incompatible joins have exactly the domain and kernel properties of
\Cref{f:lem-lower-grid}, by \Cref{f:lem-purity}. Each ordinary domain
interval has the normal-supplement and monolith properties of
\Cref{f:prop-monolith}. If at some point $i$ the nontrivial-product case
\ref{f:case-product} of \Cref{f:prop-interval-cases} occurs, it gives the
required almost simple representation of that point's link. Otherwise,
the remaining possibilities at every point are the subdirect-bottom and
the trivial-bottom full-inner cases. In both, every domain label strictly
above the base intersects the point monolith subdirectly modulo the core,
as required in \Cref{f:lem-lower-grid}. The strictness in that lemma follows
from $HV_i=T_i$: every incident domain satisfies
$T_X=H(T_X\cap V_i)$, so $T_X>H$ gives strict containment above the
point-base intersection after quotienting by $L_i$.
That lemma gives a contradiction.

For the whole lattice take the empty reverse flat as the starting top.
For a point in an initial trunk chart, its rank increment from the empty
flat is one, which lies between $1$ and $r$; its child path was therefore
reserved at the first stage. In either case at least $D-1$ generations
remain by the construction.
\end{proof}
